% Generated draft: requires mathematical and source review.
\documentclass[11pt]{article}
\usepackage[margin=25mm]{geometry}
\usepackage{fontspec}
\setmainfont{Latin Modern Roman}
\usepackage{amsmath,amssymb,mathtools,unicode-math}
\setmathfont{Latin Modern Math}
\usepackage{enumitem,needspace,xurl,hyperref}
\setlength{\parindent}{0pt}
\setlength{\parskip}{5pt}
\setlength{\emergencystretch}{3em}
\newcommand{\N}{\mathbb{N}}
\newcommand{\Z}{\mathbb{Z}}
\newcommand{\Q}{\mathbb{Q}}
\newcommand{\R}{\mathbb{R}}
\newcommand{\C}{\mathbb{C}}
\newcommand{\F}{\mathbb{F}}
\newcommand{\A}{\mathbb{A}}
\newcommand{\PP}{\mathbb{P}}
\newcommand{\E}{\mathbb{E}}
\newcommand{\eps}{\varepsilon}
\newcommand{\dd}{\,\mathrm{d}}
\begin{document}
\section*{UnsolvedMath ID: 800806}
\textbf{Source problem code:} catalogue entry
\subsection{Definitions and mathematical statement}
Let \(d\in\N\) with \(d>1\). Consider material particles moving in \(\R^d\). Write \(X_a:I\to\R^d\) for the trajectory of a particle with label \(a\), where \(I\subseteq\R\) is a time interval. The label set and the admissible class of trajectories are not specified in the supplied source. The requirement that particles stick upon collision means that, whenever the relevant trajectories are defined,
\[
X_a(t_0)=X_b(t_0)
\quad\Longrightarrow\quad
X_a(t)=X_b(t)
\qquad\text{for every }t\in I\text{ with }t\geq t_0.
\]
For a sufficiently regular Eulerian velocity field \(u\), a potential flow means a flow admitting a scalar velocity potential \(\phi\) with \(u=\nabla\phi\). The requested description must also cover non-potential flows, without imposing this gradient condition.

\textbf{Problem (Khesin, Misiolek, Modin, 2022, Problem 5).} Find a variational description of this system of sticking material particles in dimension \(d>1\), covering both potential and non-potential fluid flows. The description must be sufficiently flexible to describe the formation and subsequent development of shock waves in the system governed by the multidimensional Burgers equation.

If the intended equation is the inviscid vector Burgers equation, its smooth-region form is
\[
\partial_t u+(u\cdot\nabla)u=0,
\qquad u(t,x)\in\R^d,\quad x\in\R^d.
\]
This identification is conditional: the supplied record does not explicitly give the equation, a viscosity convention, or a notion of solution after shocks form. It also does not specify whether the particle system is discrete or continuous, its initial data and mass distribution, the velocity rule for merged particles, or the admissible variations and objective functional. Consequently, the requested variational description is an open-ended formulation problem; these missing specifications remain unresolved and are not additional hypotheses asserted here.
\subsection{Short English statement}
Find a variational formulation for particles in \(\R^d\), \(d>1\), that remain together after collision. It should accommodate both gradient and non-gradient velocity fields and describe Burgers shock formation and evolution. The source leaves the collision dynamics, solution concept, and variational framework unspecified.
\subsection{Sources}
\begin{itemize}
\item Local UnsolvedMath registry: permanent ID \(800806\), exact source problem code \texttt{catalogue entry}, title \emph{Geometric Hydrodynamics in Open Problems --- Problem 5}. The supplied record attributes the problem to Khesin, Misiolek, and Modin, with proposed year 2022. Its status is supplied as \emph{open}; this label is not an independently verified mathematical or literature assessment.
\item Khesin, Misiolek, Modin, \emph{Geometric Hydrodynamics in Open Problems} (2022), Problem 5. Supplied source link: \url{https://arxiv.org/abs/2205.01143}.
\item Supplied provenance records access on 2026-10-16 and extraction by \texttt{arxiv\_latex}. These are supplied dated provenance notes; the source and its current status have not been independently checked for this draft.
\end{itemize}
\end{document}
